For $m=2^n$ with $n\geq a$, the field $L_m$ can be replaced by
$\mathbf Q(\zeta_{d2^n})$: apply repeated-two norm compatibility
from exponent $n+a$ to $n$. This introduces no factor $2^a$.
The construction therefore covers discriminant $4h$ on the same
cofinal tail. At the coefficient level it retains the full character
$\chi_2\chi_{\rm odd}$. A dyadic character nontrivial on $\Gamma$
changes its action by the integral automorphism
$1+t\mapsto-(1+t)$, or $t\mapsto-2-t$. A dyadic sign character
instead changes the finite real-sign component. Neither operation
is a division by a sign projector.

\paragraph{5. One coefficient map produces every constituent.}
For $n$ large enough to contain the dyadic characters, put
$A_{F,n}=\operatorname{Gal}(\mathbf Q(\zeta_{F2^n})/\mathbf Q)$.
Apply Shapiro to $z_{F2^n}$, then the map on group bases
\[
 \Phi_n:\mathbf Z_2[A_{F,n}]\longrightarrow
                  \mathbf Z_2[\Gamma_n\times G],\qquad
 [s]\longmapsto\chi_0(s)[\bar s][\psi(s)].                 \tag{K10}
\]
With inverse deck actions, both sides transform under $g\in G_{\mathbf Q}$
by the inverse of its group image. On the target the extra scalar
is $\chi_0(g)$; since $\chi_0^{-1}=\chi_0$, (K10) is equivariant.
This is an integral coefficient homomorphism. Shapiro restriction
to an inertia subgroup is the direct sum of the upper restrictions;
the same zero restrictions used in (K7) therefore remain zero after
(K10). In particular the map acts on actual integral classes.
The map in (K10) is defined on this finite group for the following
reason. The constituent indexed by the trivial character of $G$ is
$\chi_0$, and
$x\circ\psi=\chi_x\chi_0$ for every $x\in\widehat G$.
All these quadratic Dirichlet characters have conductor dividing
$F2^n$ once $n$ is large enough. If an element of $G_{\mathbf Q}$
acts trivially on $\mathbf Q(\zeta_{F2^n})$, every character of $G$
is therefore one on its image under $\psi$. Character separation
makes that image the identity. Thus both $\chi_0$ and $\psi$
factor through $A_{F,n}$. Explicitly, for $g\in G_{\mathbf Q}$,
\[
 \Phi_n([g^{-1}s])
 =\chi_0(g)\,[\bar g]^{-1}[\psi(g)]^{-1}\Phi_n([s]).
\]
This verifies equivariance on each integral group basis before
evaluation at any character. The quotient to the real tower sends
each basis element to its image; its matrix has entries zero and
one. The transition maps and character evaluations are also maps
of integral coefficient modules. None of these operations is the
idempotent $|G|^{-1}\sum_g x(g)[g]$.


The maps $\Phi_n$ commute with the norm transition maps.
Proposition~8.12 has no new prime when $n$ increases, so their
images define the single class $Z_F$ in (K1).
The finite-coefficient degree-zero groups are finite and satisfy
Mittag--Leffler. The inverse-limit exact sequence
\cite[Corollary 4.1.3, p.~96]{nekovar} therefore has zero
$\varprojlim^1 H^0$ term, so it gives this degree-one class
in completed cochains.
The quotient from the full cyclotomic tower to $\Gamma$ is the
integral quotient of group bases, equivalently corestriction to
the real subfield. No factor $(1+\tau)/2$ occurs.

After evaluation at $x$, (K10) is the coefficient form of the
unnormalized corestriction weighted by $\chi_x$. Factor that
corestriction through the primitive conductor field in (K8).
Transitivity and Proposition~8.12 give the exact identity
\[
              (Z_F)_x=
       \prod_{q\mid F/d}E_q(\chi_x)\,y_{2^\infty}.        \tag{E.1}
\]
It is an identity of the given classes. The same smoothing choice
appears on both sides. There is no group-order denominator and no
comparison with a canonical zeta class of $f\otimes\chi_x$.

For clarity, at a constituent write $\chi_h=\chi_2\chi_{\rm odd}$ and
$|\Delta_h|=2^a d$. Before taking the cofinal limit, the corresponding
scalar coefficient map is exactly
\[
 [s]\longmapsto\chi_2(s)\chi_{\rm odd}(s)[s|_{\Gamma_n}],
                    \qquad n\ge\max\{a,2\}.\tag{E.2}
\]
Thus the dyadic character is part of the module action, including
when the field discriminant is $4h$.

